\section{Commercial observation systems: what a customer history must preserve}
\label{exec22:obs-scope}

A retention manager deciding whether to offer a payment concession needs a customer
history that was usable when the decision was made. A technically successful join
can still attach another person's outcome, introduce a late correction, or remove
customers whose records failed to match. These are different threats to an effect
estimate. The observation layer must therefore distinguish an account from a
person, a person from a household, and a household from a legal ownership relation.
A common telephone number can help propose a link without establishing any of
these relationships. The proposed causal library should consume evidence about
these objects, preserve alternatives, and refuse unsupported interpretations.
Commercial matching, feature and historical-storage systems provide useful parts
of that infrastructure; their documented contracts determine which parts can be
reused.

The comparison below uses AWS Entity Resolution, IBM QualityStage, Databricks
feature tables and Snowflake historical queries. It is deliberately specific about
API families and documented limits. AWS's guide is a mutable PDF generated on
22 October 2026, whose latest listed history entry dates transitive matching to
27 April 2026. The two Databricks pages report 11 September 2026 updates; IBM's
route identifies Information Server 11.7.0; the Snowflake texts do not establish an
exact service revision. These observations support a qualified technical
comparison, while leaving exact bytes at the 7 October 2026 cutoff unverified.
Informatica, Precisely, SAS data-management products and Tecton remain candidates
requiring their own technical and commercial evidence. The inspected products do
not establish an exhaustive market inventory.

\subsection{Matching outputs are hypotheses about observational units}
\label{exec22:obs-linkage}

AWS Entity Resolution separates rule-based, machine-learning and provider
workflows. The selected rule comparison distinguishes Simple and Advanced rules:
Advanced rules support fuzzy comparisons and additional Boolean expressions,
whereas Simple rules have different real-time and ID-mapping support. A generic
glossary description that calls rule matching exact-only is consequently
insufficient to describe the Advanced contract. Rule type, ordered predicates,
normalization and workflow configuration belong in the exported evidence.
The service returns a match identifier and the rule producing a match; these
identify computed groups rather than independently verified customer identities.
Its ``many-to-many'' terminology concerns comparisons across input fields sharing
a match key, such as home and mobile telephone fields. It must not be read as a
claim about account ownership.\cite{exec22_obs_aws_er_manual}

The guide also documents transitive matching as an Advanced, API-only option
configured when a new workflow is created. Records already assigned a match can
connect later records across rule levels; when multiple earlier candidates occur,
the documented resolution rule uses the smallest identifier from the earliest
rule level. Rule ordering can therefore affect grouping. The guide explicitly
warns that inappropriate order can group records with different unique-field
values. For a banking adapter, the important preserved objects are the original
rows, ordered rule specification, group membership, job identifier and provenance
of any connecting record. Treating the final group label as an immutable real-world
key would discard the evidence needed to diagnose a false merge.

The machine-learning path compares cleartext and does not support hashed input
comparisons in the selected guide. It distinguishes group-level confidence from
record-level confidence, with the latter available for incremental ML matching.
Neither the field name nor its percentage format establishes calibration for a
banking population, a posterior over complete linkage maps, or independence
between two ambiguous links. A service result should retain its native score and
scope. Converting it into a probability requires a separately justified model or
calibration study. Success and error partitions also matter: the guide describes
errors caused by invalid or repeated identifiers, and states that specified input
errors can still incur charges. Counting only the success folder misstates both
the observed population and the potential processing cost. Provider matching
additionally requires a subscription or negotiated offer, so public documentation
does not establish entitlement.\cite{exec22_obs_aws_er_manual}

IBM's versioned QualityStage overview describes a tunable probabilistic record
linkage system, a Match Designer for defining and testing specifications,
one-source and two-source matching stages, and statistics for each match pass.
These are relevant strengths for controlled specification and review. The
overview does not supply the detailed weight calculations, blocking losses,
threshold design or uncertainty exports needed to translate a score into causal
inference. Those remain prerequisites for a deeper comparison. A conditional
adapter should preserve the specification, pass and rejected records rather than
claim that the word ``probabilistic'' already supplies a coherent distribution
over customer histories.\cite{exec22_obs_ibm_matching}

Consider an exact authored example with two treated and two control records. The
treatment labels are retained in the decision system, while the outcomes arrive
through a separately linked table. Under one complete linkage map, the treated
outcomes are 8 and 6 and the control outcomes are 4 and 2. Under a second map the
two pairs are exchanged. In outcome units, the contrasts are
\begin{equation}
 \widehat\tau_{+}=7-3=4,\qquad
 \widehat\tau_{-}=3-7=-4.
 \label{exec22:obs-maps}
\end{equation}
Randomized assignment, consistency and no interference would justify interpreting
the correctly linked difference as an estimator of an average causal effect;
the four numbers alone establish only the arithmetic contrast. Let $p$ be a
stipulated probability of the first \emph{complete} map, conditional on whatever
linkage evidence is being modeled. It is not a vendor confidence score. The
map-averaged contrast is
\begin{equation}
 \overline\tau=p\widehat\tau_{+}+(1-p)\widehat\tau_{-}=8p-4.
 \label{exec22:obs-mixture}
\end{equation}
At $p=4/5$ this equals $12/5$, and at $p=2/5$ it is negative. Selecting a single
map can thus hide a decision-relevant sign ambiguity. Assigning independent
probabilities to individual links is not a substitute: it can violate exclusivity,
household structure or the requirement that a source record have a coherent
history.\cite{exec22_obs_reference}

For a separate exact uncertainty illustration, let the reported contrast equal
the selected map mean plus an independent random error taking values $-1/2$ and
$1/2$ with equal probability. Conditional variance is then $1/4$, and direct
enumeration of the four outcomes gives
\begin{equation}
 \operatorname{Var}(\widehat\tau)
 =E\{\operatorname{Var}(\widehat\tau\mid M)\}
  +\operatorname{Var}\{E(\widehat\tau\mid M)\}
 =\frac14+64p(1-p).
 \label{exec22:obs-variance}
\end{equation}
At $p=4/5$, total variance is $1049/100$, whereas a calculation conditional on one
map reports only $1/4$. This finite law demonstrates an omitted uncertainty
component; it is neither an estimated standard error for the four fixed outcomes
nor an empirical coverage result. If each additional outcome unit is stipulated
to be worth USD 10 and an offer costs USD 30, the single-map margin is USD 10,
but the map-averaged margin is
\begin{equation}
 V(p)=10(8p-4)-30=80p-70.
 \label{exec22:obs-value}
\end{equation}
It is USD $-6$ at $p=4/5$, and positive only for $p>7/8$. Actual investment value
would require valid effect estimates, calibrated linkage evidence and observed
costs. The example specifies why those inputs could change a decision.

\subsection{Event time, availability time and storage time answer different questions}
\label{exec22:obs-clocks}

Databricks documents time-series feature lookup as an AS OF join on the entity
key and a designated timestamp. Merely including a timestamp in the primary key
does not activate those semantics. The Unity Catalog path uses
\texttt{timeseries\_columns}, while the legacy Workspace Feature Store has a
different timestamp-key API. The retrieved value is the latest prior feature,
or null when no eligible value exists. A latest row with a null feature is not
silently skipped. The separate legacy-table manual describes keyed DataFrames,
Delta storage and merge behavior that leaves keys absent from the update
unchanged; it recommends the Unity Catalog successor. An adapter must pin which
client and table contract it implements.\cite{exec22_obs_databricks_pit,exec22_obs_databricks_feature}

Let $e_j$ denote the business event time of feature version $j$, $a_j$ its first
availability to the decision system, and $t$ the decision time. For a simple
immutable sequence with unique event times, a defensible historical feature is
\begin{equation}
 j^*(t)=\arg\max_{j:\ e_j\leq t,\ a_j\leq t} e_j,
 \qquad X^{\rm usable}(t)=X_{j^*(t)}.
 \label{exec22:obs-availability}
\end{equation}
An empty admissible set must yield a declared missing result. Corrections require
an additional version-selection rule among records known by $t$, rather than
allowing the latest correction to rewrite an earlier information set. If the
entity mapping itself was issued after $t$, its historical use needs a separate
restriction even when the feature record was already available.

For example, a feature value 2 has event time 8 and availability time 9; another
value 9 has event time 9 but arrives at time 12. At decision time 10, an event-time
AS OF join over today's records returns 9, while
equation~\eqref{exec22:obs-availability} returns 2. This is a constructed
counterexample to treating an event timestamp as an availability certificate,
not an observed defect in a vendor engine. A single timestamp can be sufficient
when its meaning and transformation genuinely encode the required availability
contract; that condition must be established.\cite{exec22_obs_reference}

Training and serving introduce another distinction. The Databricks manual says
that the optional lookback window constrains training and batch inference, but
online inference uses the latest value regardless of that window. In snapshot
publishing mode, the online store supports primary-key lookup without
point-in-time lookup. The proposed contract must therefore include a serving-time
freshness rule and explicit behavior when no valid history exists. Matching column
names or successfully invoking both APIs does not establish equivalent information
sets. A historical policy evaluation should use the historical lookup contract;
a live decision needs its separately checked serving contract.
\cite{exec22_obs_databricks_pit}

Snowflake Time Travel provides a complementary capability: querying, cloning and
recovering historical storage states within a retention interval. The selected
manual gives a one-day default and allows up to 90 days for permanent objects in
Enterprise Edition and above; transient and temporary objects have different
limits. Extending retention does not restore versions already outside Time Travel.
Retention also consumes storage. These are material deployment and cost
distinctions, although no price quote or commercial entitlement is established
here. A storage snapshot can preserve what a warehouse held at a particular time;
the correspondence between that state and what an operational decision-maker knew
still requires ingestion and availability evidence.\cite{exec22_obs_snowflake_time}

The SQL boundary rules must survive the adapter. \texttt{AT} includes changes at
the specified boundary; \texttt{BEFORE} a statement means immediately before its
completion, and can include another transaction committed after that statement
started. Timestamp casts, time zones and the statement-ID age limit also affect
replay. Historical queries use the current table definition, including added
columns and defaults, and current access requirements apply retroactively.
Consequently a stored query alone is insufficient to reconstruct every historical
semantic or authorization decision. Preserve the relevant schema, dictionary,
transformation, query boundary and permission evidence alongside the snapshot.
An expired historical query should remain a refusal rather than fall back to
current data.\cite{exec22_obs_snowflake_at}

\subsection{Inclusion, evaluation and a conditional integration decision}
\label{exec22:obs-selection}

Linkage failures can alter the estimand even when retained links are correct.
Suppose a randomized study has two equally common strata with causal effects 6
and 2. Let retention into the linked analysis sample be 1 in the first stratum
and $1/4$ in the second, independent of treatment and potential outcomes within
stratum. The population effect is 4, but the linked sample has stratum proportions
$4/5$ and $1/5$, giving
\begin{equation}
 \tau_{\rm linked}=\frac45\,6+\frac15\,2=\frac{26}{5}.
 \label{exec22:obs-inclusion}
\end{equation}
Weighting retained records by the inverse of their stipulated inclusion
probabilities returns 4 after normalization. This correction needs positive
inclusion probabilities and conditional transport of potential outcomes. It does
not automatically apply when matching depends on unobserved outcomes or on
post-treatment information. An audit selected only from easy, successfully linked
customers cannot establish performance in the excluded population.
\cite{exec22_obs_reference}

These distinctions support conditional reuse. AWS and QualityStage are candidates
for obtaining documented matching outputs; Databricks is a candidate for explicitly
versioned feature retrieval; Snowflake is a candidate for retained storage history.
The proposal's additional obligation is to connect those outputs to units,
information sets, target populations and causal assumptions. It currently has the
exact authored examples and design contracts described here, not qualified native
adapters. There is no basis for a blanket superiority claim over any platform.

A shared library, CLI and MCP request should declare the observational unit,
relationship type, source keys, decision time, event and availability fields,
mapping vintage, intended population and query purpose. Its result should retain
native workflow and version identifiers, alternative links or explicit absence of
them, score granularity, null/stale behavior, excluded/error records, schema and
replay evidence. Sampling uncertainty, linkage uncertainty and structural causal
ambiguity must remain distinguishable. An agent can propose schema mappings,
retrieve documentation and explain conflicting timestamps; it cannot certify
identity, ownership, causal assignment or calibration from those texts alone.
Missing availability, incompatible score semantics and expired history need typed
refusals shared by all three interfaces.

A public prototype can test these requirements without private bank data. Generate
customer histories with hidden known identities, separate household relations,
controlled name/address corruption, late corrections and known treatment effects.
Keep the truth out of the matching inputs. Compare raw exact-key joins, a declared
probabilistic matcher, corrected-clock retrieval and the proposed wrapper under
matched information. Report false merges and splits, exclusion rates by stratum,
future-information violations, effect bias, uncertainty coverage where a justified
procedure exists, and decision loss under explicit monetary scenarios. Separate
calibration from final evaluation and include abrupt population shifts. Native
vendor trials require pinned editions, rights and a reviewed budget; mathematical
examples do not substitute for them. Public prototype merit should determine
whether a later experiment with costly bank evidence is justified.
